\documentclass[10pt,a4paper]{amsart}
\usepackage{amsmath,amssymb,amsthm,graphicx}
% For commutative diagrams we use Xy-pic
\usepackage[all]{xy}

% We use 2cell for 2-commutative diagrams.
\xyoption{2cell}
% Theorem environments.
%
\theoremstyle{plain}
\newtheorem{theorem}[subsection]{Theorem}
\newtheorem{proposition}[subsection]{Proposition}
\newtheorem{lemma}[subsection]{Lemma}

\theoremstyle{definition}
\newtheorem{definition}[subsection]{Definition}
\newtheorem{example}[subsection]{Example}
\newtheorem{exercise}[subsection]{Exercise}
\newtheorem{situation}[subsection]{Situation}

\theoremstyle{remark}
\newtheorem{remark}[subsection]{Remark}
\newtheorem{remarks}[subsection]{Remarks}

\numberwithin{equation}{subsection}

% Macros
%
\def\lim{\mathop{\mathrm{lim}}\nolimits}
\def\colim{\mathop{\mathrm{colim}}\nolimits}
\def\Spec{\mathop{\mathrm{Spec}}}
\def\Hom{\mathop{\mathrm{Hom}}\nolimits}
\def\Ext{\mathop{\mathrm{Ext}}\nolimits}
\def\SheafHom{\mathop{\mathcal{H}\!\mathit{om}}\nolimits}
\def\SheafExt{\mathop{\mathcal{E}\!\mathit{xt}}\nolimits}
\def\Sch{\mathit{Sch}}
\def\Mor{\mathop{\mathrm{Mor}}\nolimits}
\def\Ob{\mathop{\mathrm{Ob}}\nolimits}
\def\Sh{\mathop{\mathit{Sh}}\nolimits}
\def\NL{\mathop{N\!L}\nolimits}
\def\CH{\mathop{\mathrm{CH}}\nolimits}
\def\proetale{{pro\text{-}\acute{e}tale}}
\def\etale{{\acute{e}tale}}
\def\QCoh{\mathit{QCoh}}
\def\Ker{\mathop{\mathrm{Ker}}}
\def\Im{\mathop{\mathrm{Im}}}
\def\Coker{\mathop{\mathrm{Coker}}}
\def\Coim{\mathop{\mathrm{Coim}}}

% Boxtimes
%
\DeclareMathSymbol{\boxtimes}{\mathbin}{AMSa}{"02}

%
% Macros for moduli stacks/spaces
%
\def\QCohstack{\mathcal{QC}\!\mathit{oh}}
\def\Cohstack{\mathcal{C}\!\mathit{oh}}
\def\Spacesstack{\mathcal{S}\!\mathit{paces}}
\def\Quotfunctor{\mathrm{Quot}}
\def\Hilbfunctor{\mathrm{Hilb}}
\def\Curvesstack{\mathcal{C}\!\mathit{urves}}
\def\Polarizedstack{\mathcal{P}\!\mathit{olarized}}
\def\Complexesstack{\mathcal{C}\!\mathit{omplexes}}
% \Pic is the operator that assigns to X its picard group, usage \Pic(X)
% \Picardstack_{X/B} denotes the Picard stack of X over B
% \Picardfunctor_{X/B} denotes the Picard functor of X over B
\def\Pic{\mathop{\mathrm{Pic}}\nolimits}
\def\Picardstack{\mathcal{P}\!\mathit{ic}}
\def\Picardfunctor{\mathrm{Pic}}
\def\Deformationcategory{\mathcal{D}\!\mathit{ef}}
\usepackage{cleveref}
\crefname{originalremark}{original remark}{original remarks}
\Crefname{originalremark}{Original remark}{Original remarks}
\makeatletter
\newlabel{remark-perf-Z-cohomology-K}{{0FEP}{source}{}{sourcepointer-1}{}}
\newlabel{remark-perf-Z-cohomology-K@cref}{{[originalremark][0][]0FEP}{[1][0][]source}{}{}{}}
\makeatother
\crefname{originalremark}{original remark}{original remarks}
\Crefname{originalremark}{Original remark}{Original remarks}
\makeatletter
\newlabel{remark-perf-Z-regular}{{0FEQ}{source}{}{sourcepointer-2}{}}
\newlabel{remark-perf-Z-regular@cref}{{[originalremark][0][]0FEQ}{[1][0][]source}{}{}{}}
\makeatother
\crefname{originallemma}{original lemma}{original lemmas}
\Crefname{originallemma}{Original lemma}{Original lemmas}
\makeatletter
\newlabel{perfect-lemma-perfect-on-regular}{{0FDC}{source}{}{sourcepointer-3}{}}
\newlabel{perfect-lemma-perfect-on-regular@cref}{{[originallemma][0][]0FDC}{[1][0][]source}{}{}{}}
\makeatother
\title{On a Noetherian regular scheme the K-theory of coherent sheaves supported on a closed subscheme agrees with the K-theory of perfect complexes supported there}
\author{The Stacks Project Authors}
\date{stacks-project:chow.tex:local-snapshot-20261012; GNU FDL 1.2 or later}
\begin{document}
\maketitle
\noindent\textit{Editorial note.} All delivered mathematical text is the author's own, copied from the pinned Stacks Project chapter file chow.tex; no statement or proof step was rewritten, reordered or removed.\par\medskip
\noindent\textit{Editorial note.} Preamble provenance: the pinned chapter file opens with an input of the shared preamble and carries no preamble text of its own, so the project's single Stacks-wide preamble staged at sources/stacks-chow/preamble.tex is excerpted verbatim as two spans, lines 14-18 (Xy-pic setup) and lines 154-217 (the native theorem declarations and mathematical macros). Nothing was added, removed or redefined; the wrapper class and the amsmath/amssymb/amsthm/graphicx base packages are supplied by the assembler.\par\medskip
\noindent\textit{Editorial note.} The counted statement and proof contain no cross reference, no citation and no included file, so the delivered body is self contained.\par\medskip
\noindent\textit{Editorial note.} Canonical location of the counted proposition: Stacks tag 0FER (sources/stacks-chow/pinned-original-source.tex line 13003).\par\medskip
\noindent\textit{Allowed standard prerequisites (editorial).}\begin{itemize}
\item On a regular scheme perfect complexes coincide with bounded complexes with coherent cohomology. Derived Categories of Schemes, Lemma labelled perfect-lemma-perfect-on-regular (Stacks tag 0FDC), invoked at proof lines 13020-13022. Supplies D\_perf(O\_X) = D\textasciicircum{}b\_Coh(O\_X), so that the cohomology sheaves of an object of D\_\{Z,perf\}(O\_X) are coherent and only finitely many are nonzero.
\item The two maps whose composite is studied, and the finite filtration of a coherent module by O\_Z-modules. Same chapter, Remarks labelled remark-perf-Z-cohomology-K (source line 12945) and remark-perf-Z-regular (source line 12978), invoked in the statement. Supply the two maps in question; the truncation argument and the filtration step are carried out in the counted span.
\end{itemize}\medskip
\ifdefined\hypertarget\hypertarget{sourcepointer-1}{}\fi
\noindent\textit{Source reference.} Original remark 0FEP: sources/stacks-chow/pinned-original-source.tex, source lines 12945--12945. Same-chapter pointer into the author environment remark at source line 12945; invoked at proof line None of the counted proof. It is retained only as a navigation pointer; every load-bearing step is inside the delivered span.\par\medskip
\ifdefined\hypertarget\hypertarget{sourcepointer-2}{}\fi
\noindent\textit{Source reference.} Original remark 0FEQ: sources/stacks-chow/pinned-original-source.tex, source lines 12978--12978. Same-chapter pointer into the author environment remark at source line 12978; invoked at proof line None of the counted proof. It is retained only as a navigation pointer; every load-bearing step is inside the delivered span.\par\medskip
\ifdefined\hypertarget\hypertarget{sourcepointer-3}{}\fi
\noindent\textit{Source reference.} Original lemma 0FDC: sources/stacks-chow/tags.txt, source lines 19212--19212. Cross-chapter standard prerequisite located in the project tag list staged in the local snapshot; invoked at proof line 13022 of the counted proof. It is retained only as a navigation pointer; every load-bearing step is inside the delivered span.\par\medskip
\begin{lemma}
\label{lemma-perf-Z-regular}
Let $X$ be a Noetherian regular scheme.
Let $Z \subset X$ be a closed subscheme. The maps constructed
in Remarks \ref{remark-perf-Z-cohomology-K} and
\ref{remark-perf-Z-regular} are mutually inverse and we get
$K'_0(Z) = K_0(D_{Z, perf}(\mathcal{O}_X))$.
\end{lemma}

\begin{proof}
Clearly the composition
$$
K_0(\textit{Coh}(Z)) \longrightarrow
K_0(D_{Z, perf}(\mathcal{O}_X)) \longrightarrow
K_0(\textit{Coh}(Z))
$$
is the identity map. Thus it suffices to show the first arrow is
surjective. Let $E$ be an object of $D_{Z, perf}(\mathcal{O}_X)$.
Recall that $D_{perf}(\mathcal{O}_X) = D^b_{\textit{Coh}}(\mathcal{O}_X)$
by Derived Categories of Schemes, Lemma
\ref{perfect-lemma-perfect-on-regular}.
Hence the cohomologies $H^i(E)$ are coherent, can be viewed
as objects of $D_{Z, perf}(\mathcal{O}_X)$, and only a finite number
are nonzero. Using the distinguished triangles of canonical truncations the
reader sees that
$$
[E] = \sum (-1)^i[H^i(E)[0]]
$$
in $K_0(D_{Z, perf}(\mathcal{O}_X))$. Then it suffices to
show that $[\mathcal{F}[0]]$ is in the image of the map
for any coherent $\mathcal{O}_X$-module set theoretically
supported on $Z$. Since we can find a finite filtration on
$\mathcal{F}$ whose subquotients are $\mathcal{O}_Z$-modules,
the proof is complete.
\end{proof}

\end{document}
